% !TEX program = xelatex
\documentclass[11pt,a4paper]{article}
\usepackage[margin=2.2cm,top=2cm]{geometry}
\usepackage{fontspec}
\usepackage{polyglossia}
\setdefaultlanguage{turkish}
\usepackage{amsmath,amssymb}
\usepackage{graphicx}
\usepackage{booktabs}
\usepackage{enumitem}
\usepackage{multicol}
\usepackage[most]{tcolorbox}
\usepackage{fancyhdr}
\usepackage{titlesec}
\usepackage{xcolor}
\usepackage{needspace}
\usepackage[hidelinks]{hyperref}

\definecolor{anamavi}{RGB}{0,45,98}
\definecolor{cozumyesil}{RGB}{0,110,50}
\definecolor{uyarikirmizi}{RGB}{180,20,20}

\titleformat{\section}{\large\bfseries\color{anamavi}}{\thesection.}{0.5em}{}
\titleformat{\subsection}{\normalsize\bfseries\color{anamavi}}{\thesubsection}{0.5em}{}

\pagestyle{fancy}\fancyhf{}
\lhead{\small MMAT105 Genel Matematik}
\rhead{\small Hafta 9 -- Ders Notu}
\cfoot{\small\thepage}
\renewcommand{\headrulewidth}{0.4pt}

\newtcolorbox{kural}[1][Kural]{colback=anamavi!5,colframe=anamavi,title=#1,fonttitle=\bfseries,breakable,left=4pt,right=4pt,top=2pt,bottom=2pt}
\newtcolorbox{dikkat}[1][Dikkat]{colback=uyarikirmizi!5,colframe=uyarikirmizi,title=#1,fonttitle=\bfseries,breakable,left=4pt,right=4pt,top=2pt,bottom=2pt}
\newcounter{ornek}
\newenvironment{ornek}{\refstepcounter{ornek}\par\medskip\noindent\textbf{\color{anamavi}Örnek \theornek.}\ }{\par}
\newenvironment{cozum}{\begin{tcolorbox}[colback=cozumyesil!6,colframe=cozumyesil,boxrule=0.5pt,left=4pt,right=4pt,top=2pt,bottom=2pt,breakable]\textbf{\color{cozumyesil}Çözüm.}\ }{\end{tcolorbox}}
\newcounter{alis}
\newenvironment{alistirma}{\refstepcounter{alis}\par\medskip\needspace{9\baselineskip}\noindent\textbf{\color{anamavi}Alıştırma \thealis.}\ }{\par}
\newcommand{\kendinizoncesi}{\needspace{12\baselineskip}}
\newcommand{\ab}[1]{\overline{#1}}
\setlength{\parindent}{0pt}
\setlength{\parskip}{3pt}


\begin{document}

% ---------------- Başlık ----------------
\begin{minipage}[c]{\textwidth}
{\LARGE\bfseries\color{anamavi} MMAT105 Genel Matematik}\\[4pt]
{\large Hafta 9 Ders Notu: Üslü Sayılar -- Özellikler,\\ Dört İşlem ve Üslü Denklemler}\\[6pt]
{\small 2026--2027 Güz}
\end{minipage}

\vspace{4pt}\noindent\textcolor{anamavi}{\rule{\textwidth}{1pt}}

\textbf{Bu notun kullanımı.} Her bölümde önce kısa bir kural kutusu, ardından çözümlü örnekler yer alır. Örnekleri çözümüne bakmadan önce kendiniz deneyin. Notun sonunda önce çözümlü alıştırmalar, ardından cevap anahtarlı ``Kendiniz Çözün'' alıştırmaları yer alır. Bu hafta üsler \textbf{tam sayıdır}; köklü sayılar ve kesirli üsler gelecek haftanın konusudur.

% =====================================================
\section{Tanım, İşaret, Sıfır ve Negatif Üs}

\begin{kural}[Tanım]
\(n\) pozitif tam sayı olmak üzere \(a^n=\underbrace{a\cdot a\cdots a}_{n\ \text{tane}}\) \ (\(a\): taban, \(n\): üs).
\begin{itemize}[leftmargin=*,itemsep=1pt]
\item \(a^1=a\); \quad \(a\neq 0\) ise \(a^0=1\); \quad \(0^0\) tanımsızdır.
\item \(1^n=1\); \quad \(n>0\) ise \(0^n=0\); \quad \((-1)^n=1\) (\(n\) çift), \ \((-1)^n=-1\) (\(n\) tek).
\item \(a>0\) ise \((-a)^{\text{çift}}>0\), \ \((-a)^{\text{tek}}<0\).
\item Negatif üs (\(a,b\neq 0\)): \quad \(a^{-n}=\dfrac{1}{a^n}\), \qquad \(\left(\dfrac{a}{b}\right)^{-n}=\left(\dfrac{b}{a}\right)^{n}\).
\end{itemize}
\end{kural}

\begin{dikkat}[Sık yapılan hatalar]
\(-a^n=-(a^n)\) ile \((-a)^n\) farklıdır: \(-3^2=-9\), \((-3)^2=9\). \ Negatif üs sayıyı negatif yapmaz: \(2^{-3}=\dfrac18\). \ Ondalık taban önce kesre çevrilir: \(0{,}2=\dfrac15\), \(0{,}25=\dfrac14\), \(1{,}5=\dfrac32\).
\end{dikkat}

\needspace{7\baselineskip}
\begin{ornek}
\(-2^4+(-2)^4-(-1)^{99}-0^{8}\) \ işleminin sonucu kaçtır?
\end{ornek}
\begin{cozum}
\(-2^4=-16\), \ \((-2)^4=16\), \ \((-1)^{99}=-1\) (\(99\) tek), \ \(0^8=0\).
\[-16+16-(-1)-0=1\]
\end{cozum}

\needspace{7\baselineskip}
\begin{ornek}
\(\left(-\dfrac13\right)^{-2}+(0{,}25)^{-1}-\left(\dfrac25\right)^{-1}\) \ işleminin sonucu kaçtır?
\end{ornek}
\begin{cozum}
\(\left(-\dfrac13\right)^{-2}=(-3)^2=9\), \quad \((0{,}25)^{-1}=\left(\dfrac14\right)^{-1}=4\), \quad \(\left(\dfrac25\right)^{-1}=\dfrac52\).
\[9+4-\frac52=\frac{26}{2}-\frac52=\frac{21}{2}\]
\end{cozum}

\needspace{7\baselineskip}
\begin{ornek}
\(2^{-1}+3^{-1}\) kaçtır? Bu sonuç \((2+3)^{-1}\)'e eşit midir?
\end{ornek}
\begin{cozum}
\(2^{-1}+3^{-1}=\dfrac12+\dfrac13=\dfrac56\). \qquad \((2+3)^{-1}=5^{-1}=\dfrac15\). \quad \(\dfrac56\neq\dfrac15\): \textbf{eşit değil}.\\
Negatif üs toplamın üzerine dağıtılamaz: \(a^{-1}+b^{-1}\neq (a+b)^{-1}\).
\end{cozum}

\needspace{7\baselineskip}
\begin{ornek}
\(a=-3\) olduğuna göre \(a^2-a^3-(-a)^2+a^0\) ifadesinin değeri kaçtır?
\end{ornek}
\begin{cozum}
\(a^2=9\); \ \(a^3=-27\Rightarrow -a^3=27\); \ \(-a=3\Rightarrow(-a)^2=9\); \ \(a^0=1\).
\[9+27-9+1=28\]
\end{cozum}

% =====================================================
\section{Üslü Sayıların Özellikleri}

\begin{kural}[Özellikler]
\(a\neq 0,\ b\neq 0\) ve \(m,n\) tam sayılar olmak üzere:\\[2pt]
\begin{tabular}{@{}ll@{}}
\(a^m\cdot a^n=a^{m+n}\), \quad \(\dfrac{a^m}{a^n}=a^{m-n}\) & (tabanlar aynı)\\[8pt]
\(a^n\cdot b^n=(ab)^n\), \quad \(\dfrac{a^n}{b^n}=\left(\dfrac{a}{b}\right)^{n}\) & (üsler aynı)\\[8pt]
\((a^m)^n=a^{m\cdot n}\) & (üssün üssü)
\end{tabular}

\textbf{Tabanları eşitleme:} \(4=2^2\), \(8=2^3\), \(16=2^4\), \(0{,}25=2^{-2}\), \(9=3^2\), \(27=3^3\), \(\frac19=3^{-2}\), \(125=5^3\), \(0{,}2=5^{-1}\), \(0{,}04=5^{-2}\).
\end{kural}

\begin{dikkat}
\(a^{m^n}=a^{(m^n)}\) ile \((a^m)^n=a^{mn}\) farklıdır: \(2^{3^2}=2^9=512\), \((2^3)^2=2^6=64\). \ \(a^m\cdot a^n\neq a^{mn}\). \ \(4^x=2^{2x}\) doğru; \(4^x=2^{x+2}\) yanlış.
\end{dikkat}

\needspace{7\baselineskip}
\begin{ornek}
\(a\neq 0\) olmak üzere \(\dfrac{(a^3)^4\cdot a^{-5}}{(a^2)^3}\) ifadesini sadeleştiriniz.
\end{ornek}
\begin{cozum}
\(\dfrac{a^{12}\cdot a^{-5}}{a^{6}}=a^{12-5-6}=a^{1}=a\).
\end{cozum}

\needspace{7\baselineskip}
\begin{ornek}
\(\dfrac{25^3\cdot 125^{-2}}{5^{-4}}\) \ işleminin sonucu kaçtır?
\end{ornek}
\begin{cozum}
\(25^3=(5^2)^3=5^6\), \ \(125^{-2}=(5^3)^{-2}=5^{-6}\). \quad \(\dfrac{5^6\cdot 5^{-6}}{5^{-4}}=5^{6-6-(-4)}=5^4=625\).
\end{cozum}

\needspace{7\baselineskip}
\begin{ornek}
\((0{,}04)^{-1}\cdot(0{,}2)^{3}\) \ işleminin sonucu kaçtır?
\end{ornek}
\begin{cozum}
\(0{,}04=5^{-2}\), \ \(0{,}2=5^{-1}\): \quad \((5^{-2})^{-1}\cdot(5^{-1})^{3}=5^{2}\cdot 5^{-3}=5^{-1}=\dfrac15=0{,}2\).\\
Kontrol: \(25\cdot 0{,}008=0{,}2\).
\end{cozum}

\needspace{7\baselineskip}
\begin{ornek}
\(\dfrac{18^5}{2^4\cdot 9^5}\) \ işleminin sonucu kaçtır?
\end{ornek}
\begin{cozum}
\(18^5=(2\cdot 9)^5=2^5\cdot 9^5\) \ \(\Rightarrow\) \ \(\dfrac{2^5\cdot 9^5}{2^4\cdot 9^5}=2^{5-4}=2\).
\end{cozum}

\needspace{7\baselineskip}
\begin{ornek}
\(10^x=a\) olduğuna göre \(0{,}01\cdot 10^{2x+1}\) ifadesini \(a\) türünden yazınız.
\end{ornek}
\begin{cozum}
\(0{,}01=10^{-2}\) \ \(\Rightarrow\) \ \(10^{-2}\cdot 10^{2x+1}=10^{2x-1}=\dfrac{(10^x)^2}{10}=\dfrac{a^2}{10}\).
\end{cozum}

% =====================================================
\section{Toplama, Çıkarma ve Bilimsel Gösterim}

\begin{kural}[Toplama -- çıkarma]
Yalnızca tabanı \textbf{ve} üssü aynı olan terimler toplanır: \(x\cdot a^n+y\cdot a^n=(x+y)\cdot a^n\).
Üsleri farklı terimlerde \textbf{en küçük kuvvet} ortak çarpan parantezine alınır: \(2^{x+2}+2^x=2^x(4+1)=5\cdot 2^x\).\\
\textbf{Bilimsel gösterim:} \(a\cdot 10^n\), \ \(1\le|a|<10\), \(n\) tam sayı. Çarpma/bölmede katsayılar ve \(10\)'un kuvvetleri ayrı ayrı işlenir.
\end{kural}

\begin{dikkat}
Toplamada üsler toplanmaz: \(2^3+2^4=24\neq 2^7\). \quad \(3^5+3^5+3^5=3\cdot 3^5=3^6\) (\(3^{15}\) değil, \(9^5\) değil).
\end{dikkat}

\needspace{7\baselineskip}
\begin{ornek}
\(\dfrac{3^{x+2}-3^{x}}{3^{x+1}+3^{x}}\) \ ifadesinin değeri kaçtır?
\end{ornek}
\begin{cozum}
Pay: \(3^x(9-1)=8\cdot 3^x\); \quad payda: \(3^x(3+1)=4\cdot 3^x\). \quad \(\dfrac{8\cdot 3^x}{4\cdot 3^x}=2\) \ (her \(x\) için).
\end{cozum}

\needspace{7\baselineskip}
\begin{ornek}
\(4^5+4^5+4^5+4^5\) \ ve \ \(2^{20}+2^{20}\) \ toplamlarını \(2\)'nin kuvveti olarak yazınız.
\end{ornek}
\begin{cozum}
\(4^5+4^5+4^5+4^5=4\cdot 4^5=4^6=(2^2)^6=2^{12}\). \qquad \(2^{20}+2^{20}=2\cdot 2^{20}=2^{21}\).
\end{cozum}

\needspace{7\baselineskip}
\begin{ornek}
\(\dfrac{5^{12}-5^{11}}{5^{11}+5^{10}}\) \ işleminin sonucu kaçtır?
\end{ornek}
\begin{cozum}
Pay: \(5^{11}(5-1)=4\cdot 5^{11}\); \quad payda: \(5^{10}(5+1)=6\cdot 5^{10}\).
\[\frac{4\cdot 5^{11}}{6\cdot 5^{10}}=\frac{4\cdot 5}{6}=\frac{20}{6}=\frac{10}{3}\]
\end{cozum}

\needspace{7\baselineskip}
\begin{ornek}
\(\dfrac{0{,}00048\cdot 25\,000\,000}{0{,}0006}\) \ işleminin sonucunu bilimsel gösterimle yazınız.
\end{ornek}
\begin{cozum}
\(\dfrac{(4{,}8\cdot 10^{-4})\cdot(2{,}5\cdot 10^{7})}{6\cdot 10^{-4}}=\dfrac{12\cdot 10^{3}}{6\cdot 10^{-4}}=2\cdot 10^{3-(-4)}=2\cdot 10^{7}\).
\end{cozum}

\needspace{7\baselineskip}
\begin{ornek}
Işık saniyede yaklaşık \(3\cdot 10^{8}\) m yol alır. Bir yıl yaklaşık \(3{,}15\cdot 10^{7}\) saniyedir. Işığın bir yılda aldığı yol (ışık yılı) yaklaşık kaç metredir?
\end{ornek}
\begin{cozum}
\((3\cdot 10^{8})\cdot(3{,}15\cdot 10^{7})=(3\cdot 3{,}15)\cdot 10^{8+7}=9{,}45\cdot 10^{15}\) m.
\end{cozum}

% =====================================================
\section{Sıralama}

\begin{kural}[Sıralama]
\(a>1\) ise üssü büyük olan büyüktür; \(0<a<1\) ise üs büyüdükçe sayı küçülür. Üsler eşitse (pozitif tabanlarda) tabanı büyük olan büyüktür. Farklı taban ve üslerde ya tabanlar eşitlenir ya da üsler, üslerin EBOB'u ile eşitlenir.
\end{kural}

\needspace{7\baselineskip}
\begin{ornek}
\(a=(0{,}5)^{-10}\), \ \(b=4^{6}\), \ \(c=8^{3}\) \ sayılarını küçükten büyüğe sıralayınız.
\end{ornek}
\begin{cozum}
\(a=(2^{-1})^{-10}=2^{10}\), \ \(b=2^{12}\), \ \(c=2^{9}\). \ Taban \(2>1\): \(9<10<12\) \(\Rightarrow\) \(c<a<b\).
\end{cozum}

\needspace{7\baselineskip}
\begin{ornek}
\(a=2^{45}\), \ \(b=3^{30}\), \ \(c=7^{15}\) \ sayılarını küçükten büyüğe sıralayınız.
\end{ornek}
\begin{cozum}
Üslerin EBOB'u \(15\): \ \(a=(2^3)^{15}=8^{15}\), \ \(b=(3^2)^{15}=9^{15}\), \ \(c=7^{15}\). \ \(7<8<9\) \(\Rightarrow\) \(c<a<b\).
\end{cozum}

% =====================================================
\section{Üslü Denklemler}

\begin{kural}[Üslü denklemler]
\begin{itemize}[leftmargin=*,itemsep=1pt]
\item \(a\neq 0,\ a\neq 1,\ a\neq -1\) ise: \ \(a^x=a^y\Rightarrow x=y\) \quad ve \quad \(a^x=1\Rightarrow x=0\).
\item \(a>0,\ b>0,\ a\neq b\) ise: \ \(a^x=b^x\Rightarrow x=0\).
\item Taban bilinmeyen \(x^a=x^b\) (\(x\) tam sayı): üsler eşit olabilir; ya da taban \(1\) olabilir; taban \(0\) olabilir (iki üs de pozitifse); taban \(-1\) olabilir (iki üs de çift ya da ikisi de tekse). Her aday kontrol edilir.
\end{itemize}
\end{kural}

\begin{dikkat}
Üsleri farklı terimlerin toplamı varsa önce ortak çarpan parantezi: \(2^{x+3}+2^{x}=9\cdot 2^x\). \ \(1^x=1^y\) ve \((-1)^2=(-1)^4\) olduğundan \(a=0,\pm 1\) tabanlarında üsler eşitlenemez.
\end{dikkat}

\needspace{7\baselineskip}
\begin{ornek}
\(27^{x-1}=9^{x+1}\) \ denklemini çözünüz.
\end{ornek}
\begin{cozum}
\(3^{3x-3}=3^{2x+2}\) \(\Rightarrow\) \(3x-3=2x+2\) \(\Rightarrow\) \(x=5\). \quad Kontrol: \(27^4=3^{12}=9^6\).
\end{cozum}

\needspace{7\baselineskip}
\begin{ornek}
\((0{,}2)^{x+1}=125^{x-3}\) \ denklemini çözünüz.
\end{ornek}
\begin{cozum}
\((0{,}2)^{x+1}=(5^{-1})^{x+1}=5^{-x-1}\), \ \(125^{x-3}=5^{3x-9}\).
\[-x-1=3x-9\ \Rightarrow\ 8=4x\ \Rightarrow\ x=2\qquad(\text{Kontrol: } 0{,}2^3=0{,}008=125^{-1})\]
\end{cozum}

\needspace{7\baselineskip}
\begin{ornek}
\(2^{x+3}+2^{x+1}-2^{x}=72\) \ denklemini çözünüz.
\end{ornek}
\begin{cozum}
\(2^x(8+2-1)=9\cdot 2^x=72\) \(\Rightarrow\) \(2^x=8=2^3\) \(\Rightarrow\) \(x=3\). \quad Kontrol: \(64+16-8=72\).
\end{cozum}

\needspace{7\baselineskip}
\begin{ornek}
\(2^{x-3}=3^{x-3}\) \ ve \ \(5^{y+1}=1\) \ olduğuna göre \(x+y\) kaçtır?
\end{ornek}
\begin{cozum}
Tabanlar farklı, üsler aynı: \(x-3=0\Rightarrow x=3\). \quad \(5^{y+1}=5^0\Rightarrow y=-1\). \quad \(x+y=2\).
\end{cozum}

\needspace{7\baselineskip}
\begin{ornek}
\(x\) tam sayı olmak üzere \((x+1)^{2x}=(x+1)^{x+3}\) \ denkleminin çözüm kümesi nedir?
\end{ornek}
\begin{cozum}
\begin{itemize}[leftmargin=*,itemsep=0pt,topsep=0pt]
\item Üsler eşit: \(2x=x+3\Rightarrow x=3\) (taban \(4\)) \checkmark
\item Taban \(1\): \(x=0\): \(1^0=1^3=1\) \checkmark
\item Taban \(0\): \(x=-1\): \(0^{-2}\) tanımsız \(\times\)
\item Taban \(-1\): \(x=-2\): \((-1)^{-4}=1\), \((-1)^{1}=-1\) \(\times\)
\end{itemize}
Çözüm kümesi: \(\{0,\ 3\}\).
\end{cozum}

\needspace{7\baselineskip}
\begin{ornek}\textit{(Birler basamağı hatırlatma)}
\(8^{43}\) sayısının birler basamağındaki rakam kaçtır?
\end{ornek}
\begin{cozum}
\(8\)'in kuvvetlerinin birler basamakları \(8,4,2,6\) biçiminde \(4\)'lü tekrar eder. \(43=4\cdot 10+3\) \(\Rightarrow\) üçüncü terim: \(2\).
\end{cozum}

% =====================================================
\section{Özet: Formül Kartı}
\begin{center}
\renewcommand{\arraystretch}{1.6}
\begin{tabular}{@{}lp{12.9cm}@{}}
\toprule
Konu & Formül / Kural\\\midrule
Tanım & \(a^n=a\cdot a\cdots a\) (\(n\) tane); \ \(a^1=a\); \ \(a^0=1\) (\(a\neq 0\)); \ \(0^0\) tanımsız\\
İşaret & \((-a)^{\text{çift}}>0\), \((-a)^{\text{tek}}<0\) (\(a>0\)); \ \(-a^n\neq(-a)^n\)\\
Negatif üs & \(a^{-n}=\dfrac{1}{a^n}\), \ \(\left(\dfrac ab\right)^{-n}=\left(\dfrac ba\right)^n\)\\
Çarpma -- bölme & \(a^m a^n=a^{m+n}\), \ \(\dfrac{a^m}{a^n}=a^{m-n}\), \ \(a^nb^n=(ab)^n\), \ \(\dfrac{a^n}{b^n}=\left(\dfrac ab\right)^n\)\\
Üssün üssü & \((a^m)^n=a^{mn}\neq a^{m^n}\)\\
Toplama & benzer terimler: \(xa^n+ya^n=(x+y)a^n\); \ ortak çarpan: \(2^{x+2}+2^x=5\cdot 2^x\)\\
Bilimsel gösterim & \(a\cdot 10^n\), \ \(1\le|a|<10\)\\
Sıralama & taban \(>1\): büyük üs büyük; \ \(0<\)taban\(<1\): büyük üs küçük; \ ya tabanlar ya üsler (EBOB ile) eşitlenir\\
Denklem & \(a^x=a^y\Rightarrow x=y\), \ \(a^x=1\Rightarrow x=0\) (\(a\neq 0,\pm1\)); \ \(a^x=b^x\Rightarrow x=0\) (\(a,b>0,\ a\neq b\))\\
\(x^a=x^b\) & üsler eşit; taban \(1\); taban \(0\) (üsler pozitif); taban \(-1\) (üsler aynı paritede)\\\bottomrule
\end{tabular}
\end{center}

% =====================================================
\newpage
% Hafta 9 -- alıştırma gövdesi (ders notu ve ayrı alıştırma kağıdı tarafından \input edilir)
\section{Çözümlü Alıştırmalar}
Her soruyu önce kendiniz deneyin, sonra çözümü okuyun.

\begin{alistirma}
\((-2)^3+(-1)^{2023}-(-4)^2+3^0\) \ işleminin sonucu kaçtır?
\end{alistirma}
\begin{cozum}
\((-2)^3=-8\) (tek kuvvet), \ \((-1)^{2023}=-1\) (\(2023\) tek), \ \((-4)^2=16\) (çift kuvvet), \ \(3^0=1\).\\
Sonuç: \(-8+(-1)-16+1=-24\).
\end{cozum}

\begin{alistirma}
\(\left(\dfrac23\right)^{-2}+\left(\dfrac12\right)^{-3}\) \ işleminin sonucu kaçtır?
\end{alistirma}
\begin{cozum}
Negatif üste taban ters çevrilir: \(\left(\dfrac23\right)^{-2}=\left(\dfrac32\right)^{2}=\dfrac94\), \quad \(\left(\dfrac12\right)^{-3}=2^3=8\).
\[\frac94+8=\frac94+\frac{32}{4}=\frac{41}{4}\]
\end{cozum}

\begin{alistirma}
\((0{,}2)^{-1}+(0{,}5)^{-2}\) \ işleminin sonucu kaçtır?
\end{alistirma}
\begin{cozum}
\(0{,}2=\frac15\), \ \(0{,}5=\frac12\) \(\Rightarrow\) \(\left(\frac15\right)^{-1}+\left(\frac12\right)^{-2}=5+2^2=5+4=9\).
\end{cozum}

\begin{alistirma}
\((2^3)^4\cdot(2^{-2})^5\cdot 8^2\) \ işleminin sonucu kaçtır?
\end{alistirma}
\begin{cozum}
Üssün üssü: \((2^3)^4=2^{12}\), \ \((2^{-2})^5=2^{-10}\), \ \(8^2=(2^3)^2=2^6\).\\
\(2^{12}\cdot 2^{-10}\cdot 2^{6}=2^{12-10+6}=2^{8}=256\).
\end{cozum}

\begin{alistirma}
\(\dfrac{3^{11}+3^{10}}{3^{10}-3^{9}}\) \ işleminin sonucu kaçtır?
\end{alistirma}
\begin{cozum}
Pay kısmında en küçük kuvvet \(3^{10}\), paydada \(3^{9}\) ortak çarpana alınır:
\[\frac{3^{10}(3+1)}{3^{9}(3-1)}=\frac{4\cdot 3^{10}}{2\cdot 3^{9}}=2\cdot 3=6\]
\end{cozum}

\begin{alistirma}
\(8^{x+1}=16^{x-1}\) \ denklemini çözünüz.
\end{alistirma}
\begin{cozum}
İki taraf \(2\) tabanında yazılır: \(8^{x+1}=2^{3(x+1)}=2^{3x+3}\), \ \(16^{x-1}=2^{4(x-1)}=2^{4x-4}\).
\[3x+3=4x-4\ \Rightarrow\ x=7\qquad(\text{Kontrol: } 8^8=2^{24}=16^6)\]
\end{cozum}

\begin{alistirma}
\(2^{x+1}+2^{x}+2^{x-1}=56\) \ denklemini çözünüz.
\end{alistirma}
\begin{cozum}
En küçük kuvvet \(2^{x-1}\) parantezine alınır: \(2^{x+1}=4\cdot 2^{x-1}\), \ \(2^{x}=2\cdot 2^{x-1}\).
\[2^{x-1}(4+2+1)=7\cdot 2^{x-1}=56\ \Rightarrow\ 2^{x-1}=8\ \Rightarrow\ x=4\qquad(\text{Kontrol: } 32+16+8=56)\]
\end{cozum}

\begin{alistirma}
\(a=2^{60}\), \ \(b=3^{40}\), \ \(c=5^{20}\) \ sayılarını küçükten büyüğe sıralayınız.
\end{alistirma}
\begin{cozum}
Üslerin EBOB'u \(20\): \ \(a=(2^3)^{20}=8^{20}\), \ \(b=(3^2)^{20}=9^{20}\), \ \(c=5^{20}\).\\
Üsler eşit olduğundan tabanlar karşılaştırılır: \(5<8<9\) \ \(\Rightarrow\) \ \(c<a<b\).
\end{cozum}

\kendinizoncesi
\section{Alıştırmalar (Kendiniz Çözün)}

\begin{enumerate}[label=\textbf{\arabic*.},leftmargin=*,itemsep=6pt]
\item \(-5^2+(-5)^2-(-2)^4\) \ işleminin sonucu kaçtır?
\item \(2^{-3}+3^{-1}-4^{-1}\) \ işleminin sonucu kaçtır?
\item \((0{,}1)^{-2}+(0{,}5)^{-1}\) \ işleminin sonucu kaçtır?
\item \(2^x\cdot 2^5\cdot 2^{-3}\) \ ifadesini tek bir üslü sayı olarak yazınız.
\item \(\dfrac{3^{12}\cdot 3^{-4}}{3^5}\) \ işleminin sonucu kaçtır?
\item \(9^3\cdot 27^{-2}\cdot 81\) \ işleminin sonucu kaçtır?
\item \(14\,500\,000\) \ ve \ \(0{,}000025\) \ sayılarını bilimsel gösterimle yazınız.
\item \((6\cdot 10^{-3})\cdot(5\cdot 10^{8})\) \ işleminin sonucunu bilimsel gösterimle yazınız.
\item \(2^x=a\) olduğuna göre \(8^{x+1}\) ifadesini \(a\) türünden yazınız.
\item \(3^{2x-1}=243\) \ denklemini çözünüz.
\item \(\left(\dfrac13\right)^{x-1}=9^{x+2}\) \ denklemini çözünüz.
\item \(3^{x+2}-3^{x}=72\) \ denklemini çözünüz.
\item \(a=4^{12}\), \ \(b=3^{18}\), \ \(c=2^{30}\) \ sayılarını küçükten büyüğe sıralayınız.
\item \(x\) tam sayı olmak üzere \((x-1)^{2x}=(x-1)^{x+4}\) \ denklemini sağlayan kaç farklı \(x\) değeri vardır?
\end{enumerate}

\subsection*{Cevap Anahtarı}
{\small
\begin{multicols}{2}
\begin{enumerate}[label=\textbf{\arabic*.},leftmargin=*,itemsep=1pt]
\item \(-25+25-16=-16\).
\item \(\frac18+\frac13-\frac14=\frac{3+8-6}{24}=\frac{5}{24}\).
\item \(10^2+2=102\).
\item \(2^{x+5-3}=2^{x+2}\).
\item \(3^{12-4-5}=3^3=27\).
\item \(3^{6}\cdot 3^{-6}\cdot 3^{4}=3^4=81\).
\item \(1{,}45\cdot 10^{7}\); \ \(2{,}5\cdot 10^{-5}\).
\item \(30\cdot 10^{5}=3\cdot 10^{6}\).
\item \(8\cdot 8^x=8\cdot(2^x)^3=8a^3\).
\item \(3^{2x-1}=3^5\Rightarrow 2x-1=5\Rightarrow x=3\).
\item \(3^{-x+1}=3^{2x+4}\Rightarrow -x+1=2x+4\Rightarrow x=-1\).
\item \(3^x(9-1)=72\Rightarrow 3^x=9\Rightarrow x=2\).
\item \(a=2^{24}\); EBOB \(6\): \(16^6<27^6<32^6\Rightarrow a<b<c\).
\item \(x=4\) (üsler eşit), \(x=2\) (taban \(1\)), \(x=1\) (\(0^2=0^5\)), \(x=0\) (\((-1)^0=(-1)^4\)): \(4\) değer.
\end{enumerate}
\end{multicols}}


\vspace{6pt}
\noindent\textcolor{anamavi}{\rule{\textwidth}{0.5pt}}\\
{\small Gelecek hafta: Köklü sayılar.}

\end{document}
